%-----------------------------------------------------------
% Modularity and Adaptability
% \label{subsec:modularity}
%-----------------------------------------------------------

A central principle in the phantom's design is modularity — the ability to replace, adapt, or extend 
individual components without the need to re-manufacture the entire assembly. 
This approach supports multiple experimental configurations, facilitates rapid prototyping of new geometries, 
and makes the system reusable by other researchers who may wish to test alternative antenna layouts 
or conductive materials.

Examples include:
\begin{itemize}
    \item \textbf{Base connector:} a parametric connector geometry that can be extended or reoriented 
          to reach any desired location within the gel volume. 
          This enables the creation of new antenna types or electrode placements while keeping the rest 
          of the model unchanged.
    \item \textbf{Lid:} located at the lower end of the shell (neck region), it serves as the main interface 
          for filling the internal cavity with the conductive gel and for fixing or routing the antenna connectors. 
          The lid can be reprinted with customized openings to accommodate different experimental setups, 
          while a \textbf{base (uncut) version} is provided as a reusable template for future modifications.
\end{itemize}

This modular structure simplifies experimentation and ensures that both hardware and design files 
can be easily adapted or shared, promoting reproducibility and collaborative development.

\begin{figure}[h]
    \centering
    \includegraphics[width=0.65\textwidth]{figures/base_connector.png}
    \caption{Parametric base connector model extending towards the internal conductive medium.}
    \label{fig:base_connector}
\end{figure}

\nota{agregar imgs desp.}
